% !TEX encoding = UTF-8 Unicode
\documentclass[11pt, a4paper, twoside, openright]{book}

% Setup of packages used
\usepackage[onehalfspacing]{setspace}		% One half spacing
\usepackage[a4paper, 					% Page Layout
                     showframe,					% This shows the frame
                     twoside, includehead,
                     footskip=7mm, headsep=6mm,
                     top=25mm, bottom=25mm, inner=30mm, outer=25mm]{geometry}
%\usepackage{sansmathfonts}				% Sans Serif equations
\renewcommand*\familydefault{\sfdefault} 		% Sans Serif as default font
\usepackage[utf8]{inputenc}				% Encoding of files: utf8
\usepackage[T1]{fontenc}					% Output font encoding for international characters

\usepackage{titlesec}					% Redefine chapter and cahpter* titles
\titleformat{\chapter}[display]{\huge}{\vspace{-1.5cm}\hfill \textbf{\chaptertitlename\ \thechapter}}{-5pt}{\huge \hfill}[%
{\titlerule[1.2pt] \enlargethispage{-0.75\baselineskip}}]
\titleformat{name=\chapter,numberless}[display]{\huge}{\vspace{-1.5cm}\vspace{-1ex}}{-5pt}{\huge \hfill}[%
{\titlerule[1.2pt] \enlargethispage{-0.75\baselineskip}}]

\usepackage{fancyhdr}					% Redefine headers
\pagestyle{fancy}
\fancyhf{}								% No header nor footer
\fancyhead[LE,RO]{\thepage}				% Left even and right odd Page Number
\fancyhead[RE]{\textit{\leftmark}}			% Right even Section
\fancyhead[LO]{\textit{\rightmark}}			% Left odd Chapter

\fancypagestyle{plain}{					% To change the footer on chapter and chapter*
	\fancyhf{}							% No header nor footer
	\fancyfoot[C]{\vspace{-11mm}\thepage}	% Footer with number displaced
	\renewcommand{\headrulewidth}{0pt}	% No line on header
	\renewcommand{\footrulewidth}{0pt}		% No line on footer
}

\RequirePackage{caption} 				% Caption customization
\captionsetup{justification=centerlast,font=small,labelfont=sc,margin=1cm}

\usepackage{hyperref}					% Hyperlinks on pdf
\hypersetup{
    colorlinks=true,
    linkcolor=blue,
    filecolor=magenta,      
    urlcolor=blue,
    citecolor=blue,    
}

% Setup the language and its properties
\usepackage[spanish, es-tabla, es-nodecimaldot]{babel}

\usepackage{graphicx}
%\graphicspath{ {figs/} }					% Use this if custom figures folder is needed

\usepackage[square, numbers]{natbib}		% Bibliography
\usepackage{tikz}						% Required for title page
\usetikzlibrary{babel}						% Required to make TikZ works with babel
\usepackage[section]{placeins}				% To flush floats before new section
%\nobibliography*						% To make bibliography works

% User packages, packages used in this document
\usepackage{lipsum}
\usepackage{layouts}
\printinunitsof{mm}

% Macros provided
\def\fecha{\ifcase\month\or
  Enero\or Febrero\or Marzo\or Abril\or Mayo\or Junio\or
  Julio\or Agosto\or Septiembre\or Octubre\or Noviembre\or Diciembre\fi \space\number\year
}

% User Definitions
\title{Diseño, Modelado y Control de un Convertidor Multinivel Fuente de Corriente}
\author{Esp. Ing. Pablo \textsc{Cossutta}}
\date{\fecha}

\frontmatter
\begin{document}
\begin{titlepage}
	\onehalfspacing
	\enlargethispage{0.65\baselineskip}
	\begin{tikzpicture}[remember picture, overlay]
		\coordinate (top_right) at 
		    ([xshift=-2.5cm, yshift=-2.5cm]current page.north east);
		\coordinate (top_left) at 
		    ([xshift=2.3cm, yshift=-1.8cm]current page.north west);
		\coordinate (bottom_right) at 
		    ([xshift=-1.8cm, yshift=1.8cm]current page.south east);
		\node[inner sep=0, anchor=north east] at (top_right) {\href{http://www.itba.edu.ar}{\includegraphics[height=19mm, trim={180 200 200 200}, clip]{figs/logo_itba.png}}};
		\draw[double, line width = 0.5pt] (top_left) rectangle (bottom_right);
	\end{tikzpicture}
	\par
%	\begin{large}
		\vspace{-0.1cm}
		\noindent \textbf{DEPARTAMENTO DE INVESTIGACIÓN Y DOCTORADO}\par
		\vspace{3cm}
		\begin{center}
			{\Huge \textbf{\@title}\par}
			{\huge \textbf{Subtítulo del trabajo (si lo hubiere)}\par}
		\end{center}
		\vfill
		\noindent \textbf{AUTOR:} \@author \par
		\vspace{1cm}
		\noindent \textbf{DIRECTOR:} Dra. María Inés \textsc{Valla}\par
		\noindent \textbf{CO-DIRECTOR:} Dr. Miguel \textsc{Aguirre}\par	
		\vspace{1cm}
		\noindent{TESIS PRESENTADA PARA LA OBTENCIÓN DEL TÍTULO DE}\par
		\noindent\textbf{DOCTOR EN INGENIERÍA}\par
		\vspace{1cm}
		\noindent\textbf{Comité evaluador}\par
		\begin{normalsize}		
			\noindent%
			\makebox[0.33\textwidth][l]{Dr./Dra. Nombre \textsc{Apellido}}
			\makebox[0.33\textwidth][c]{Dr./Dra. Nombre \textsc{Apellido}} 
			\makebox[0.33\textwidth][r]{Dr./Dra. Nombre \textsc{Apellido}}\par
		\end{normalsize}
		\vspace{1cm}
		\begin{center}
			\textbf{CIUDAD AUTÓNOMA DE BUENOS AIRES}\\
			\textbf{\@date}\par
		\end{center}
%	\end{large}
\end{titlepage}

\mainmatter
\pagenumbering{arabic}
\chapter{Primero}
\pagevalues
\section{Lorem Ipsum}
\lipsum[1-15]
\begin{figure}
	\pagediagram
	\caption{Diagrama de página}
	\label{fig1}
\end{figure}
La Fig. \ref{fig1} muestra el formato que adopta Latex para el Layout.
\end{document}  
